\question{4.m4}{Van de kansvariabele $X$ is de verdelingsfunctie als volgt:}
    \begin{center}
        \begin{tabular}{cccccccc}
            \toprule
                {\bfseries Uitkomst $k$} & $10$ & $11$ & $12$ & $13$ & $14$ & $15$ & $16$\\
            \cmidrule{1-1} \cmidrule{2-2} \cmidrule{3-3} \cmidrule{4-4} \cmidrule{5-5} \cmidrule{6-6} \cmidrule{7-7} \cmidrule{8-8}
                $F(k)$ & $0.16$ & $0.28$ & $0.46$ & $0.66$ & $0.84$ & $0.94$ & $1.00$\\
            \bottomrule
        \end{tabular}
    \end{center}
    Bereken $P(X=14)$. Wat is de uitkomst?

\begin{enumerate}[label=(\alph*)]
    \item $0.18$
    \item $0.84$
    \item $0.10$
    \item $0.16$
\end{enumerate}
\answer{
    In dit geval is niet de kansfunctie $f(k)$ gegeven, maar de (cumulatieve) verdelingsfunctie $F(k) = P(X \leq k)$, oftewel de CDF.
    De kans dat $X$ de waarde $14$ aanneemt, oftewel $P(X=14)$, kunnen we berekenen door het verschil te nemen tussen de verdelingsfunctie bij $14$ en de verdelingsfunctie bij de vorige uitkomst, namelijk $13$:
    \[
        P(X=14) = P(X\leq 14) - P(X\leq 13) = F(14) - F(13) = 0.84 - 0.66 = 0.18.
    \]
    Het juiste antwoord is (a).
}